\section{Background}
\label{sec:background}

\subsection{A hierarchical hybrid}

CMT~\citep{guo2022cmt} keeps the four-stage pyramid of convolutional
backbones and puts attention inside it. A stem of three $3\times3$
convolutions, the first strided, reduces the input before any token exists;
each of the four stages then begins with a strided convolution that halves
the spatial resolution and widens the channels, so the token count falls by
four from stage to stage. In the configuration we use the four stages hold
$2$, $2$, $10$, and $2$ blocks.

A block has three parts. A local perception unit adds a depthwise $3\times3$
convolution with a residual connection, giving each token access to its
spatial neighbours without attention. A lightweight multi-head self-attention
follows, in which keys and values are first spatially reduced by a depthwise
strided convolution---by factors of $8$, $4$, $2$, and $1$ in the four
stages---so that early stages, where tokens are numerous, do not pay the full
quadratic cost. A learned relative position bias is added to the attention
logits. The block closes with an inverted-residual feed-forward network: a
pointwise expansion, a depthwise $3\times3$ convolution, and a pointwise
projection.

One consequence of the relative position bias matters later. Its table is
sized by the token grid and by the reduced grid it attends to, so it grows
with the square of the token count rather than staying fixed, which makes the
model's parameter count a function of input resolution.

\subsection{A flat state space backbone}

Vision Mamba~\citep{zhu2024vim} keeps the flat layout of a plain vision
transformer---one patch embedding, one resolution throughout, learned
absolute position embeddings, a class token placed in the middle of the
sequence---and replaces self-attention with a selective state space
scan~\citep{gu2023mamba}.

The scan carries a hidden state of width $N$ along the sequence, updating it
at each token by a discretised linear recurrence. Its continuous parameters
are discretised by zero-order hold with a step size that is itself produced
from the input, so the state transition is input dependent rather than fixed:
the model can decide, per token, how much of the state to retain. Around the
scan the block has an input projection that expands the channel width, a
short depthwise convolution along the sequence, the projections that produce
the input-dependent parameters, and an output projection. In the
configuration we use the state width is $16$ and the scan is run in both
directions over the sequence with separate weights, which is what makes the
backbone usable for images, where nothing distinguishes one end of the
sequence from the other.

Because the layout is flat, an image reaches the scan as a row-major
flattening of the token grid. Two horizontally adjacent tokens are
neighbours in the sequence; two vertically adjacent tokens are separated by
the width of the grid. Later hierarchical state space
backbones~\citep{liu2024vmamba,huang2024localmamba} address this with
multi-directional or windowed scans, and we return in
Section~\ref{sec:factorial} to why our factorial design deliberately does
not.

As the flat attention counterpart we use DeiT~\citep{touvron2021deit}, which
is the baseline the state space backbone's own authors compare against.

\subsection{Where the six claims come from}

The claims we examine are of two kinds, and the difference decides what a
failure to reproduce means.

Three are values published by the original authors. The cost of a block is
compared analytically: self-attention needs $2M^2d$ operations for the
attention matrices and about $4Md^2$ for its projections, while a scan needs
$8MDN$ per direction, linear in the token count $M$. Setting the quadratic
term against the linear ones puts a cross-over near $M \approx 2d$, and a
table of block costs at a large token count reports a ratio of more than an
order of magnitude past it. Separately, the state space backbone's authors
report a throughput advantage of $2.8\times$ and a memory reduction of
$86.8\%$ at an input resolution of $1248^2$ on their hardware. And the
hybrid's own paper reports an ablation table attributing its accuracy to the
hierarchy, the convolutional stem, the local perception unit, and the
inverted-residual network.

Three are assertions about mechanism, made in comparison writing rather than
measured: that the scan direction leaves an anisotropic imprint on the
effective receptive field; that a model without softmax normalisation escapes
the dilution of attention weight over a large object and therefore favours
large objects; and that the hybrid's component-wise ablation explains the
accuracy gap between the two families, which requires the ablation to
transfer across a change of token-mixing operator. Two of the three appear in
our own earlier version of this comparison.

Quoting a reported value into new conditions and asserting a mechanism
without measuring it fail differently. We treat the first as a limit on the
range over which a reported result has been shown to hold, and the second as
a claim we are entitled to test directly.
